\documentclass[11pt]{article}
\usepackage[a4paper,margin=2.5cm]{geometry}
\usepackage[T1]{fontenc}
\usepackage{charter}
\usepackage{booktabs}
\usepackage{graphicx}
\usepackage{xcolor}
\usepackage[hidelinks]{hyperref}
\usepackage{enumitem}
\setlength{\parskip}{4pt}
\setlength{\parindent}{0pt}
\definecolor{boxgrey}{gray}{0.95}
\newenvironment{example}{\begin{quote}\small\begin{sffamily}}{\end{sffamily}\end{quote}}

\title{\Large\bfseries Keeping the Future Out of the Past:\\
The Book Filter, Every Step of the Way}
\author{Ranke project --- decontamination pipeline validation}
\date{24 July 2026 --- Version 3.0 (adds Step 7, the upgrade round, and the production blueprint)}

\begin{document}
\maketitle

\begin{small}
\noindent\textbf{Document header.} Work performed 22--24 July 2026 on the office workstation (RTX 4090) plus
two brief rented cloud sessions. Code and manifests: \texttt{github.com/huggingbeard/ranke-edition-dating},
folder \texttt{experiments/content-dating/}. Data, judgments and reports:
\texttt{D:\textbackslash UZHechist Dropbox\textbackslash Joachim Voth\textbackslash claudecode\textbackslash
ranke-contamination-audit\textbackslash experiments\textbackslash content-dating\textbackslash}.
Day-by-day trail: that project's \texttt{LOG.md}. Section~\ref{sec:docs} maps every detailed document.
All parameter choices remain provisional pending the professor's review.
\end{small}

\tableofcontents

\section{The mission}
We train language models --- the Ranke models --- that are supposed to know the world only up to a
cutoff year, 1929. They are used as ``synthetic respondents'' from the past, so a single leak of later
knowledge poisons the well: a model that has read about the 1960s cannot honestly tell us what 1929 felt
like. The training library of roughly two million scanned books must therefore contain no text
\emph{written} after 1929.

Note the word \emph{written}. A book printed in 1963 that faithfully re-sets an 1850 novel contains
1850s words --- valuable, not poisonous (though its modern introduction, if it has one, is poison).
Conversely, a book whose catalogue card says ``1865'' can turn out to contain a 1913 congressional
record. Paperwork lies; text does not. Everything below follows from that asymmetry.

Our guiding rule throughout, set at the start of the project: \textbf{throwing away a clean old book is
cheap (we have millions); keeping one contaminated book is expensive.} So we always prefer catching more
contamination even at the price of discarding more clean books.

\section{The starting point: the front-pages check and its blind spot}
The production filter (the ``front-pages check'') reads only about 3\% of each book: the title page,
the copyright page, and the catalogue entry. Its rule is simple: take the newest year found in the
catalogue, in an earlier automated content reading, or printed near a publishing word
(``copyright,'' ``printed,'' ``edition''), and discard the book if that year is later than 1929.
A guard ignores years that sit next to library-stamp words (``received,'' ``date due,'' ``digitized''),
because a library receiving a book in 1940 says nothing about when it was written.

This check is excellent when the paperwork is honest --- and blind when it is not. In an earlier audit
of 35{,}572 kept books, exactly three modern books had slipped through. All three had the same disguise:
wrong-or-old catalogue years \emph{and} no readable publication date in their scanned front pages. One
of them opens like this (verbatim from the scan):

\begin{example}
``\ldots Harvard College Library FROM THE BRIGHT LEGACY One half the income from this Legacy, which was
received in 1880 under the will of JONATHAN BROWN BRIGHT of Waltham, Massachusetts, is to be expended
for books for the College Library\ldots''
\end{example}

A bookplate about an 1880 bequest, catalogue year 1909 --- and further inside, the text of a 1913
congressional record. Nothing on the front pages betrays it. The question this week: can a second
check, one that \emph{reads the body of the book}, close that blind spot --- and what does the whole
two-check system then achieve?

\section{Step 1 --- building an exam with an answer key}
You cannot grade a filter without books whose true dates you know. Assembling that answer key was
half the work, and it went wrong once in an instructive way.

\subsection{Clean old books (the ``false-alarm'' side)}
We drew 999 books that five independent date signals all said were pre-1929, then had two different AI
examiners (unrelated to the readers being tested later) each read up to 24 windows of 3{,}000
characters spread through each book, asking: does anything here require post-1929 knowledge?
\textbf{918 books passed both examiners} and became our certified clean set. The 81 the stricter
examiner worried about were set aside in a review queue --- spot checks suggest over-caution (one
``modern'' flag was the bracketed reference number ``(1981, 3026)'' in a citation), but we did not
argue with the protocol: uncertain books simply do not get certified. Cost of this certification:
\$4.21 of API calls.

\subsection{Modern books --- and the trap the professor caught}
Our first batch of known-modern books came from the same scanned-library collection, filtered to
publication year 1947 or later. The catch: in a public-domain collection, post-1947 books are almost
all \emph{government documents} --- legislature reports, hearings, agency annual reports. A filter
scored against those is, in the professor's words, ``a gov-doc classifier by another name'': it can
pass the test by recognising bureaucratic prose rather than by dating content. Measured later, this
easy set flattered every reader by 10--25 points (Figure~\ref{fig:strata}).

The fix was two harder sets. First, \textbf{600 ordinary authored books from 1930--1946} --- poetry,
fiction, local histories, textbooks; books whose language is barely distinguishable from the 1920s.
(Finding them meant scanning the catalogue metadata of all 983{,}004 books in the collection, title-screening
1{,}816 candidates, and having both examiners certify each one as genuinely newly-written --- 631 of
899 examined passed.) Second, \textbf{322 gold-audited contaminants from our own corpus}: books a
careful audit had confirmed as post-1929 editions, including 239 with modern text, 80 modern
\emph{reprints} of old works (kept separate --- their body text is genuinely old), and the three
escapees themselves.

\section{Step 2 --- the reading test}
\subsection{The question we ask}
Every reader in the test received book excerpts with exactly this instruction (verbatim):

\begin{example}
``Read this passage from a scanned book and judge: was its text WRITTEN before 1929, or after 1929?
Old, foreign-language, poetic, or historical writing is BEFORE. Modern prose, a post-1929
event/person/technology, or a modern perspective is AFTER. Ignore OCR noise. Passage: [\ldots]
Answer with only one word: BEFORE or AFTER.''
\end{example}

Each book contributed up to 16 randomly-placed excerpts (plus 16 aimed at suspicious vocabulary) at
four lengths from 250 to 2{,}000 characters --- 232{,}680 excerpts in all. Ten AI readers of different
sizes judged every one of them overnight on the office computer; that let us ask, for every reader and
every rule (``how many excerpts, how long, how many AFTER votes to flag a book?''), what it catches
and what it falsely flags --- 648 configurations per reader.

\subsection{The always-say-modern trap}
The first results table contained a booby trap worth remembering. The smallest reader appeared to be
the champion: it ``caught'' 99.6\% of modern books from a single excerpt. It also flagged
\textbf{99.9\% of the certified-old books}. It had simply learned to answer AFTER to everything ---
perfect recall, useless filter. \emph{Every} catch-rate we quote is therefore paired with its
false-alarm rate; a number without that pair is meaningless.

\begin{figure}[h]\centering
\includegraphics[width=0.86\textwidth]{fig_cloud.png}
\caption{\small\textbf{Figure 1.} Every dot is one reader in one configuration. The horizontal axis is
the share of 918 certified-old books wrongly flagged (false alarms); the vertical axis is the share of
239 known modern-content books caught. Good filters live in the top-left. The grey dashed line is the
0.95 catch-rate target; the dotted line marks 10\% false alarms. The top-left corner stays empty: no
single reader reaches the target at an acceptable false-alarm rate. The grey cloud at top-right is the
small reader that says ``modern'' to everything.}
\label{fig:cloud}
\end{figure}

\subsection{What the best single readers achieve --- and why there is a ceiling}
At 10\% false alarms or better, the best single reader (a 20-billion-parameter model, ``gpt-oss'')
catches 85 of 100 known contaminants, but only 73 of 100 of the harder 1930s ordinary books. We read
the missed books by hand. They are not model failures so much as an impossibility: their words carry
no date. Three verbatim examples of body text from books the best reader missed --- each one
\emph{certainly} post-1929 (clear copyright evidence, audit confidence at least 0.95):

\begin{example}
(devotional poetry, copyright 1947) ``\ldots thy shrine; How sweetly dost thou show thy face; When thou
dost own me thine; O Father, let me praise thy name In great humility\ldots''\\[3pt]
(county history, copyright 1940, quoting an 1811 newspaper) ``\ldots The steamboat New Orleans came to
anchor before this port on Monday last. This is a large and handsome vessel, burthen 350 tons\ldots''\\[3pt]
(taxonomy monograph, printed 1947) ``\ldots I found that this was at best only an average difference and
that no definite line between \emph{O.~pulchella} and \emph{O.~concinna} could be drawn\ldots''
\end{example}

A 1947 hymn in King-James diction; an old newspaper quoted verbatim; timeless Latin taxonomy. No
body-reader, human or machine, can date these from the words. But notice: every one of these misses
carries its date \emph{on its front pages} --- which is exactly what the first check reads. The two
checks fail on different books. That observation drives everything that follows.

\begin{figure}[h]\centering
\includegraphics[width=0.86\textwidth]{fig_strata.png}
\caption{\small\textbf{Figure 2.} Catch rates of the four competent readers at each one's best
setting with false alarms held under 10\%. Blue: the easy government-document set --- it flatters every
reader. Red: gold-audited contaminants. Purple: 1930s ordinary books, the hardest and most decision-relevant
set. The dashed line is the 0.95 target; nothing reaches it on the purple bars.}
\label{fig:strata}
\end{figure}

\section{Step 3 --- readers working in teams, and the big reader}
\subsection{Agreement crushes false alarms}
All ten readers judged the \emph{same} excerpts, so team rules cost nothing to test. The best one:
a passage counts as ``modern'' only when \textbf{two readers independently agree on that same passage}.
Each reader hallucinates ``AFTER'' on different passages, so demanding agreement collapses false
alarms --- which in turn lets each reader run at a more aggressive, higher-catch setting. The 20B model
paired with a 14B model this way: 90 of 100 contaminants and 80 of 100 hard 1930s books caught at
9\% false alarms --- five to seven points better than any single reader at the same cost in clean books.
We tested 12{,}528 team combinations in all, and (to avoid crowning a lucky winner) picked teams on one
half of the books and report their scores only from the other half.

\subsection{The 35-billion-parameter reader}
The professor pressed for one more contender, a 35B model. Its story in brief: it does not fit the
office card at standard compression (two attempts froze the machine); at the professor's suggestion
(``more quantization''), a 3-bit version ran correctly but 25 times slower than the 20B. To settle its
true ability we rented a large cloud card for half an hour (cost: \$0.74, machine verified shut down
three times afterwards) and ran the \emph{uncompressed} model. Verdict: it is the sharpest reader
tested --- alone it catches 89/87 (contaminants / 1930s books) at 8.6\% false alarms, and teamed with
the 20B it reaches \textbf{91 and 88 of 100 at under 10\% false alarms} --- the best honest operating
point of the whole exercise. (Pleasingly, the 3-bit version scores within 1--3 points of the
uncompressed one: compression was never the problem, only speed.)

\section{Step 4 --- testing the whole pipeline, not the parts}
Adding the two checks' rates would be wrong twice over: the second check only ever sees books the
first check \emph{passed}, and the two fail on correlated books (no-imprint books are often also
archaic-sounding books --- the escapee class). So we ran books through both checks \emph{in sequence}.

\subsection{Disguised books, verified realistic}
The crucial test stratum: modern books wearing the escapees' disguise. We took 155 certified-modern
books and corrupted their paperwork the way reality does --- cutting away the scanned front pages at a
natural page boundary (scans genuinely often start mid-book) and, in 40\% of cases, pasting in a
\emph{genuine} old library bookplate harvested from our certified-clean books. The professor's
condition was that disguises must be indistinguishable from real scans, so we tested exactly that: a
blind AI judge shown shuffled openings called 80\% of our disguises ``artificially assembled'' --- and
85\% of \emph{real} book scans too (old scans just look strange). Telling them apart, it scored 47.5\%,
i.e.\ coin-flip: the disguises pass. (Also reassuring: 39\% of genuine old books have no readable
publication date either --- our disguises hide in a large, real crowd.)

\subsection{What the sequence achieves}
\begin{itemize}[leftmargin=1.4em,itemsep=2pt]
\item \textbf{Honest paperwork: the first check alone catches 99.6--100\%.} The body-reading check adds
      nothing on these --- as predicted from the miss examples above.
\item \textbf{Disguised books: the first check catches none} (that is what the disguise does), and the
      body-reading team rescues the great majority --- 93 of 100 in the tuning data, later corrected to
      88 of 100 out of sample (Step 5).
\item \textbf{One subtlety we could not test directly:} one input of the first check is a year estimate
      produced by a 14B model during the original corpus build, and its author is on holiday. Instead
      of guessing, we graded that model's existing homework: on 101 audited books with wrong-old
      paperwork where we know the true year, it read the \emph{true} year 93 times --- and its estimate
      tracked the truth, not the (wrong) paperwork, in 88\% of cases, so wrong catalogue cards do not
      mislead it. Folding that in: roughly \textbf{99 of 100 disguised books are removed overall}, and
      the survivors are precisely the archaic-content, no-imprint, wrong-paperwork triple --- the
      escapee class, of which body-reading catches about one in three.
\item \textbf{Reprints, by the professor's rule:} flagged reprints get one more reading focused on
      editorial matter --- ``any later-date introduction or editor's notes? then out; a clean re-set
      stays.'' Of 80 reprints, 50 carried modern apparatus (e.g.\ verbatim findings:
      ``Introduction by S.~Barnitz Williams (1958)''; ``Preface dated February, 1952'') and 48 of those
      reached removal; \textbf{27 clean re-sets were rescued} --- premium period text a blanket ban
      would have burned.
\item \textbf{The price:} 13 of 100 certified-old books are wrongly discarded (a stricter voting
      setting halves this, at a known cost in catch rate).
\end{itemize}

\section{Step 5 --- the honesty test: predictions first, fresh books second}
Numbers measured on the books used to tune a system flatter it. So before drawing a single new book we
wrote five predictions --- each with an acceptable range --- into the public project repository
(time-stamped commit \texttt{5176fd8}). Then we assembled \textbf{743 fresh books that no earlier step
had ever touched} (365 newly certified old books, 250 newly certified honest moderns, 128 newly
disguised books built with the frozen recipe), ran the frozen filter, and compared. Full statistics
--- exact confidence intervals and significance tests --- are in \texttt{score\_oos.py} and
\texttt{oos\_results.csv}; the verdict table:

\begin{center}\small
\begin{tabular}{@{}llll@{}}
\toprule
\textbf{Prediction (made in advance)} & \textbf{Allowed range} & \textbf{Fresh-book result} & \textbf{Verdict}\\
\midrule
Honest moderns caught (composite) & at least 0.98 & 1.000 (250 of 250) & pass\\
Body-check adds $\le$0.02 on honest moderns & $\le$0.02 & +0.000 & pass\\
Disguised books rescued by body check & 0.86--0.97 & 0.875 (112 of 128) & pass$^{*}$\\
Disguised books removed overall & 0.98--1.00 & 0.991 & pass\\
Clean books wrongly discarded & 0.06--0.14 & 0.132 (48 of 365) & pass (at the edge)\\
\bottomrule
\end{tabular}
\end{center}

{\small $^{*}$Within the pre-registered range, but measurably below the tuning-data figure of 0.93
(the gap is statistically real, $p=0.023$). This is the expected ``winner's optimism'' of tuning data,
and it is why the honest numbers quoted throughout this document are the fresh-book ones: rescue
\textbf{0.88}, clean-book loss \textbf{0.13}.}

All five predictions held. Nothing was re-tuned on the fresh books.

\section{Step 6 --- a third check: read everything, cheaply}
The professor asked a final question: the body-readers only \emph{sample} a book (about 16{,}000 of
300{,}000 characters) --- could a brute-force scan of \emph{every} character, looking for modern
vocabulary and dated years, catch what the samplers miss? The idea has history: it is exactly how the
three original escapees were found.

Building it taught a lesson worth recording. Our first version flagged 60\% of \emph{certified-old}
books, for two embarrassing reasons: the term-matcher found ``nato'' \emph{inside} Italian words like
``impersonato,'' and the AI adjudicator, shown excerpts labelled ``[marker:nato],'' dutifully declared
they mentioned NATO --- trusting our label over the text, which contained no such thing. The fixed
version matches whole words only, accepts years only next to publishing words (``copyright,''
``published,'' ``printed''), shows the adjudicator \emph{unlabelled} raw excerpts, and requires it to
quote the exact proving phrase. That collapsed false flags on old books from 60\% to 1--3\% while
keeping 28--50\% flag rates on the books the readers had missed.

Measured end to end, the third check rescues about \textbf{30 of 100} reader-missed contaminants,
\textbf{27 of 100} missed 1930s books, and \textbf{50 of 100} missed fresh-test moderns --- but
\textbf{0 of 16} disguised books (the disguise removes the very pages its year-scan reads). It costs
about half a clean book per 100 (and one of its four ``false'' alarms turned out to be a plausible
genuine discovery: a supposedly-clean book containing ``converted into chambers in the year 1963'').
Verdict: worth adding --- nearly free, +3--5 points on the hard sets, and an occasional detective ---
but it does not break the 0.95 ceiling either. Adjudicator failure modes, for the record: a 1900s
engineer projecting ``the year 2000'' (innocent), an OCR-garbled ``1999,'' and one world-knowledge slip.

\section{Step 7 --- the upgrade round: a bigger reader, and better search terms}
Two final ideas from the professor, both tested the same day.

\textbf{A bigger reader? No.} We rented the large cloud card again and ran the 120-billion-parameter
sibling of our 20B reader over every test excerpt (cost \$6.69, including a re-run after a third of
its answers were cut off mid-thought --- lesson: thinking models need room to finish). Verdict: it
ties the 35B on known contaminants but does \emph{worse} --- by five points --- on the hardest 1930s
books, while running four times slower on identical hardware. Bigger stopped helping; the 35B's
judgement, not raw size, is what matters. We can stop shopping for readers.

\textbf{Better search terms? Yes, with a twist.} The old suspicious-vocabulary list was built for a
1945 cutoff, so it ignores 1930s language entirely. We built a 1929-cutoff list two ways: by hand
(history: ``Luftwaffe,'' ``penicillin,'' ``nylon''\ldots) and by mining --- asking which words appear
in at least five of our known-modern books and in \emph{zero} of 1,361 certified old books. The mining
found something better than events: the everyday fingerprints of modern English. ``Basically.''
``Contacted.'' ``Phoned.'' ``Sponsored.'' Nobody wrote those words before 1929. (The zero-hit test
also vetoed three of our hand-picked terms --- ``new deal,'' ``nazi,'' ``stalin'' all occur innocently
in pre-1929 books --- machines checking humans, usefully.) Net effect: the third check's rescue rate
on missed 1930s books rose from 27 to 39 in 100. One honest limit: a judge required to quote
\emph{proof} rightly refuses to condemn a book for containing the word ``basically,'' so these
statistical fingerprints screen well but convict rarely; turning them into a formal voting rule is a
designed-but-untested option.

\section{The production blueprint}
Run on rented cloud cards, the full two-million-book job looks like this: the free paperwork rule
first (plus the reprint reading); the two-reader team over the survivors; the exhaustive term-scan
with targeted adjudication over what remains. Money and time, from measured speeds:
about \textbf{160 rented-GPU hours $\approx$ \$250} (conservative; $\approx$\$150 with standard
batching optimisations), finishing in \textbf{about two days on four cards} or one day on eight ---
against roughly \$9{,}000 for the na\"ive one-big-model approach. Full stage-by-stage specification,
with every frozen parameter and dial: \texttt{PIPELINE.md} in the repository.

\section{The bottom line}
\begin{itemize}[leftmargin=1.4em,itemsep=2pt]
\item Books with honest paperwork: \textbf{essentially 100\%} removed by the front-pages check alone.
\item Books with wrong or missing paperwork --- the realistic threat: \textbf{about 99 of 100} removed
      by the full sequence; the residual is the small class whose words and paperwork are both silent.
\item With the third check added, reader-level catch rates on the hardest sets rise a further 3--5
      points (1930s books 0.87\,$\to$\,0.90; audited contaminants 0.90\,$\to$\,0.93) for about half a
      clean book per 100 extra; the three checks fail on provably \emph{different} books, which is the
      whole design.
\item Reprints: modern apparatus removed, clean re-sets rescued.
\item Cost: about \textbf{13 clean books lost per 100} at the chosen setting (dial available); money
      and compute in the table below.
\item Open items: the 81-book review queue (looks like examiner over-caution; awaiting adjudication);
      archiving the original 14B year-estimator's code when its author returns; and the final choice
      between the standard and strict voting settings.
\end{itemize}

\subsection*{What the whole exercise cost}
\begin{center}\small
\begin{tabular}{@{}p{8.2cm}rr@{}}
\toprule
\textbf{Item} & \textbf{Time} & \textbf{Dollars}\\
\midrule
API calls, measured via the billing meter (certifying 999+400 old books \$4.21+\$1.9; certifying
1{,}499 modern candidates \$1.07+; title screens, realism test, examiner spot-checks) & --- & \$7.21\\
Rented cloud A100 card, five sessions (35B verification; fresh-test judging; 120B evaluation incl.\
its re-run), each terminated and verified off & $\sim$5.5 h & \$7.81\\
Office RTX 4090: overnight reader sweep ($\sim$10.5 h), team/delta passes ($\sim$2.5 h), 35B local
attempts ($\sim$2.5 h), reprint + fresh-test passes ($\sim$0.5 h) & $\sim$16 h & $\sim$\$1.50 electricity\\
\midrule
\textbf{Total} & \textbf{$\sim$18 GPU-hours local + 5.5 cloud-hours} & \textbf{$\sim$\$17}\\
\bottomrule
\end{tabular}
\end{center}
{\small For scale: dating all two million books through the full two-check pipeline would cost roughly
\$75--110 (first check, mostly free rules), plus either $\sim$21 days of the office 4090 or $\sim$\$180
of rented A100 time (body-reading check, depending on the reader team chosen). The same job done
na\"ively with a large commercial model would run to roughly \$9{,}000.}

\section{Where the details live}\label{sec:docs}
\begin{small}
\begin{tabular}{@{}p{5.9cm}p{9.3cm}@{}}
\toprule
\textbf{Document} & \textbf{What it contains} \\
\midrule
\texttt{PIPELINE.md} & \textbf{The authoritative production spec}: every stage, frozen parameter, dial, and cost line. \\
\texttt{content\_dating\_report.html} (v2.2) & Full technical report: all readers, all configurations, figures, verbatim misses, production section. \\
\texttt{PREREGISTRATION\_OOS.md} & The five predictions and rules, committed before the fresh test (\texttt{5176fd8}). \\
\texttt{score\_oos.py}, \texttt{oos\_results.csv}, \texttt{oos\_plants\_results.csv} & Fresh-test scoring code and per-book results with confidence intervals. \\
\texttt{results.csv} / \texttt{ensemble\_results.csv} / \texttt{combo35\_results.csv} & All 648 single-reader, 2{,}016 team, and 12{,}528 combined configurations. \\
\texttt{composite\_results.csv} & Per-book two-check outcomes for every panel stratum. \\
\texttt{reprint\_verdicts.csv} & The 80 reprint decisions with quoted evidence. \\
\texttt{manifest\_*.csv} (8 files) & Exactly which books are in every stratum (IDs, years, labels). \\
\texttt{build\_plants.py} / \texttt{test\_plant\_realism.py} & The disguise recipe and its blind realism test. \\
\texttt{review\_queue\_negatives.csv} & The 81 examiner-flagged old books awaiting adjudication. \\
\texttt{notes\textbackslash 2026-07-23-pipeline-composite-eval-design.md} & The evaluation design memo. \\
\texttt{LOG.md} (project root) & The complete trail, including everything that broke and how it was fixed. \\
\bottomrule
\end{tabular}
\end{small}

\end{document}
